% !TeX root = ../../Book2.tex
\chapter{变分法}
\label{b2:ch:14}
\BookChapterNumber{b2:ch:14}
% 本章摘要（不计入编号）
\begin{chapterabstract}
本章用函数空间中的微分研究积分泛函，推导标量与向量值 Euler--Lagrange 方程，并从边界项辨认自然边界条件。Dirichlet 能量说明凸性怎样把驻点变成唯一极小点；波动作用量则说明驻点未必是极小点。最后把第一卷的直接方法应用于梯度幂能量，逐项验证强制性、弱闭约束与弱下半连续性，得到在一般有界开集上的存在性，并说明弱方程与经典正则性之间仍有距离。
\end{chapterabstract}

% 本章目标（不计入编号）
\begin{learninggoals}
\begin{itemize}
\item\label{b2:item:14-goal-derivative}区分方向变分、Gâteaux 导数与 Fréchet 导数，检查积分下求导所需的控制。
\item\label{b2:item:14-goal-boundary}从允许扰动出发推导内部方程及自然边界条件，核对外法向的符号。
\item\label{b2:item:14-goal-energy}用精确能量差或第二变分判断极小性，区别作用量与守恒能量。
\item\label{b2:item:14-goal-direct}在 Sobolev 边界类中实施直接方法，指出强制性、弱紧性和下半连续性各自的作用。
\end{itemize}
\end{learninggoals}

一条连接固定端点的曲线怎样使长度最小？一张固定在边框上的薄膜怎样达到平衡？这类问题比较的是整个函数，而不只是有限个参数。微积分仍然有效，但导数、紧性和边界条件都必须放到合适的函数空间中。为避免把光滑形式计算误当作存在性证明，本章先推导候选者必须满足的方程，再独立构造极小元。

\ProseParagraphBreak

本章使用第一卷第九章的弱收敛与凸泛函理论、第二十二章的 Sobolev 空间，以及第二十四章的 Poincaré 不等式；具体调用处均写明结论。关于变分法的系统叙述，可进一步阅读 Bernard Dacorogna 的《Direct Methods in the Calculus of Variations》\cite{Dacorogna2008Direct}及 Lawrence C. Evans 的《Partial Differential Equations》\cite{Evans2010Partial}。本章只使用下面完整核验的凸能量模型，不预借更一般的非凸积分泛函下半连续性理论。

% !TeX root = ../../Book2.tex
\section{分离变量法}
\label{b2:sec:1201}

有限杆的温度既由初温决定，也由端点如何交换热量决定。空间模态要先满足端点约束，其时间衰减再由相应特征值决定。本节把分离变量所得级数放进已经建立的能量空间，随后用镜像表示及具体边界数据检验它的含义。

\subsection{有界区间}
\label{b2:subsec:120101}

热传导中，温度梯度引起从高温向低温的通量。设细杆的热容量密度为 \(w(x)>0\)、传导系数为 \(p(x)>0\)，通量取 \(-p(x)u_x\)。无热源时，对任意子区间的热量平衡给出 \(wu_t=(pu_x)_x\)。若再有与温度成正比的损耗 \(q(x)u\)，则方程成为
\[
 wu_t=(pu_x)_x-qu,
 \quad u_t+Au=0.
\]
在有限区间及\cref{b2:def:10-sl-operator}的连续系数、分离边界假设下，空间算子就是\chapref{b2:ch:10}建立的 Sturm--Liouville 算子。本节先取非负能量情形；例如 \(q\geq0\)，以及 Robin 系数非负。

\ProseParagraphBreak

对初温 \(f\in H=L^2(I,w\,\mathrm{d}x)\)，\cref{b2:prop:10-heat-energy-limit}给出
\begin{equation}\label{b2:eq:12-interval-expansion}
 u(t)=\sum_ne^{-\lambda_nt}f_ne_n,
 \quad f_n=\langle f,e_n\rangle_H.
\end{equation}
它在零时以 \(H\) 范数取得初值，正时间属于 \(D(A)\)，并在 \(H\) 中满足方程；还在 \(L^2((0,T);V)\) 中具有有限能量。连续系数并不使空间解自动无限次可微：当 \(p\) 只有连续性时，正则性应以 \(pu_x\) 和 \(D(A)\) 表述。下面转到常系数杆，才进一步证明所有空间阶数的经典光滑性。

\subsection{Fourier 级数}
\label{b2:subsec:120102}

在 \(0<x<\ell\) 上固定 \(\kappa>0\)，考虑 \(u_t=\kappa u_{xx}\)。齐次 Dirichlet 条件给出
\begin{equation}\label{b2:eq:12-dirichlet-series}
 u_D(t,x)=\sum_{n=1}^\infty f_ne^{-\kappa(n\pi/\ell)^2t}
 \sqrt{\dfrac2\ell}\sin\dfrac{n\pi x}{\ell}.
\end{equation}
齐次 Neumann 条件给出
\begin{equation}\label{b2:eq:12-neumann-series}
 u_N(t,x)=\overline f+
 \sum_{n=1}^\infty g_ne^{-\kappa(n\pi/\ell)^2t}
 \sqrt{\dfrac2\ell}\cos\dfrac{n\pi x}{\ell},
 \quad\overline f=\dfrac1\ell\int_0^\ell f.
\end{equation}
其中 \(f_n,g_n\) 分别为正弦和非零余弦系数。常数项来自零模态；它不随时间衰减。

\Needspace{6\baselineskip}
\begin{proposition}[正时间的经典解]\label{b2:prop:12-interval-smooth}
设 \(\ell,\kappa>0\)，\(u_D,u_N\) 分别为零 Dirichlet 和零 Neumann 边界的热模态级数\eqref{b2:eq:12-dirichlet-series}、\eqref{b2:eq:12-neumann-series}，系数取初值的标准正交正弦、余弦系数。若 \(f\in L^2((0,\ell);\mathbb C)\)，上面两个级数及其任意有限阶 \(x,t\) 导数在 \([\delta,T]\times[0,\ell]\) 上绝对一致收敛，其中 \(0<\delta<T\)。因此它们在正时间光滑，满足各自方程和边界条件。
\end{proposition}
\begin{proof}
求 \(j\) 次空间导数及 \(m\) 次时间导数只会把第 \(n\) 项乘以至多 \(C_{j,m}n^{j+2m}\)。Cauchy--Schwarz 给出统一可和上界
\[
 \sum_{n\geq1}|f_n|n^{j+2m}e^{-c\delta n^2}
 \leq\left(\sum_{n\geq1}|f_n|^2\right)^{1/2}
 \left(\sum_{n\geq1}n^{2j+4m}e^{-2c\delta n^2}\right)^{1/2}<\infty.
\]
余弦系数的证明相同。每个阶数都有相应一致收敛，故逐项微分合法。有限项已经满足方程和边界条件，极限保持它们。
\end{proof}

例如两端温度固定为零，初温在杆内恒为 \(1\)，则
\[
 u_D(t,x)=\dfrac4\pi\sum_{j=0}^{\infty}
 \dfrac{e^{-\kappa((2j+1)\pi/\ell)^2t}}{2j+1}
 \sin\dfrac{(2j+1)\pi x}{\ell}.
\]
正时间的端点值为零，而零时的 \(L^2\) 初温仍是常数一。它不可能在闭区间上随着 \(t\downarrow0\) 一致趋于一，因为端点值始终为零；初边值在角点不相容并不妨碍平方能量解存在。

\subsection{边界条件}
\label{b2:subsec:120103}

周期热核提供另一种表示。沿用第二编的实线核 \(H_s(x)=(4\pi s)^{-1/2}e^{-x^2/(4s)}\)，定义
\begin{align*}
 K_D(t,x,y)&=\sum_{m\in\mathbb Z}
 \left(H_{\kappa t}(x-y+2m\ell)-H_{\kappa t}(x+y+2m\ell)\right),\\
 K_N(t,x,y)&=\sum_{m\in\mathbb Z}
 \left(H_{\kappa t}(x-y+2m\ell)+H_{\kappa t}(x+y+2m\ell)\right).
\end{align*}
则 \(u_D(t,x)=\int_0^\ell K_D(t,x,y)\cdot f(y)\,\mathrm{d}y\)，Neumann 情形用 \(K_N\)。两式分别来自 \(f\) 的奇延拓、偶延拓及 \(2\ell\) 周期化；将卷积在每个周期分割后换元即得。高斯尾保证对正时间紧区间的一切求和与微分可交换。第二编的周期 Fourier 系数唯一性证明它们与式\eqref{b2:eq:12-dirichlet-series}、式\eqref{b2:eq:12-neumann-series}一致。

\ProseParagraphBreak

这两个空间表达式在小时间有直接用途：附近的镜像点控制边界影响，远处项指数小。大时间则应使用特征级数，只保留低频项。由于奇延拓与偶延拓在有限 \(p\) 范数中有界，周期热核逼近还给出 \(f\in L^p(0,\ell)\)、\(1\leq p<\infty\) 时的初值收敛。Dirichlet 的闭区间一致收敛要求 \(f\) 连续且两端为零；Neumann 情形对任意连续 \(f\) 都成立，零时不要求导数边界相容。

\ProseParagraphBreak

非齐次端点先用提升处理。例如 \(u(t,0)=a(t)\)、\(u(t,\ell)=b(t)\)，其中 \(a,b\in C^1\)，令
\[
 r(t,x)=a(t)\left(1-\dfrac x\ell\right)+b(t)\dfrac x\ell,
 \quad v=u-r.
\]
那么 \(v\) 的边界为零，且 \(v_t-\kappa v_{xx}=-r_t\)。记源项正弦系数为 \(F_n(t)\)，逐模态公式为
\[
 v_n(t)=e^{-\kappa\lambda_nt}v_n(0)
 +\int_0^te^{-\kappa\lambda_n(t-s)}F_n(s)\,\mathrm{d}s.
\]
对 \(F\in C([0,T];L^2)\)，积分项的平方范数由 \(\int_0^t\|F(s)\|_2\,\mathrm{d}s\) 控制，有限截断因紧时间区间上的连续像为紧集而一致逼近。因而该公式给出 \(C([0,T];L^2)\) 中的解；不能只因源项连续就额外断言所有高阶导数连续。

\ProseParagraphBreak

非齐次端温还可给出一个完全显式的周期响应。在半直线 \(x>0\) 上取端温 \(u(t,0)=\cos(\omega t)\)，其中 \(\omega>0\)。令 \(\alpha=\sqrt{\omega/(2\kappa)}\)，则
\[
 u(t,x)=e^{-\alpha x}\cos(\omega t-\alpha x)
\]
满足热方程与这个边界条件。验证时先用复值函数 \(e^{\mathrm{i}\omega t-(1+\mathrm{i})\alpha x}\)：其时间导数乘以 \(\mathrm{i}\omega\)，二阶空间导数乘以 \((1+\mathrm{i})^2\alpha^2=\mathrm{i}\omega/\kappa\)，再取实部即可。它的初温是 \(e^{-\alpha x}\cos(\alpha x)\)，因此此式描述已经建立的周期响应；若改给另一份初温，还须加上零端温差解，不能直接沿用同一式子。距离增加既压低振幅，也使相位推迟。振幅降为端点的 \(e^{-1}\) 倍所需距离为 \(\alpha^{-1}=\sqrt{2\kappa/\omega}\)，所以快速变化的端温主要影响靠近端点的一段区域。

\ProseParagraphBreak

Robin 条件的解仍用式\eqref{b2:eq:12-interval-expansion}，但总热量还受内部损耗影响。为将边界通量写得有意义，取有限区间 \(I=(a,b)\)、实连续系数 \(p,w,q\in C([a,b])\)、\(p,w>0\)。在开时间区间 \(J\subset(0,\infty)\) 上，设
\[
 u\in C^1(J;L^2(I,w\,\mathrm{d}x))\cap C(J;H^1(I)),
 \qquad pu_x\in C(J;H^1(I)),
\]
并在 \(L^2(I)\) 中满足 \(wu_t=(pu_x)_x-qu\)。此时通量 \(pu_x\) 有绝对连续代表及两个端点的迹，而积分是 \(L^2(I,w\,\mathrm{d}x)\) 上的连续线性泛函。因此可以先对时间求导，再用空间基本定理，得到
\begin{equation}\label{b2:eq:12-weighted-heat-balance}
 \frac{\mathrm{d}}{\mathrm{d}t}\int_a^b w(x)u(t,x)\,\mathrm{d}x
 =[p(x)u_x(t,x)]_a^b-\int_a^b q(x)u(t,x)\,\mathrm{d}x,
 \qquad t\in J.
\end{equation}
前面的谱解在任意闭正时间带上具有这些强解性质：指数衰减控制 \(Au\) 及 \(u_t\) 的 \(H\) 范数连续性。由\cref{b2:lem:10-form-coercive}及 \((pv')'=qv-wAv\)，还有
\[
 \|v\|_{H^1(I)}+\|pv'\|_{H^1(I)}
 \leq C\bigl(\|v\|_H+\|Av\|_H\bigr),\qquad v\in D(A).
\]
其中第一项先由移位能量控制，通量的零阶范数再由 \(v'\) 控制，一阶范数由所示方程控制。因此 \(u\) 与通量也具有刚才要求的 \(H^1\) 时间连续性。这里不需要把仅连续的 \(p\) 作经典微分，也不要求一般初温在零时具有边界导数。

\ProseParagraphBreak

式\eqref{b2:eq:12-weighted-heat-balance}的第一项是边界净输入，第二项是内部吸收或产生。Neumann 条件只消除前一项；即使边界绝热，内部反应也会改变总热量。例如在 \((0,\ell)\) 上取 \(p=w=q=1\)、齐次 Neumann 条件及初温 \(f=1\)，则 \(u(t,x)=e^{-t}\)，加权总热量为 \(\ell e^{-t}\)，并不守恒。

对无内部反应的纯热方程，取 \(w=1\)、\(p=\kappa>0\)、\(q=0\)，才得到
\[
 \frac{\mathrm{d}}{\mathrm{d}t}\int_0^\ell u(t,x)\,\mathrm{d}x
 =\kappa\bigl(u_x(t,\ell)-u_x(t,0)\bigr).
\]
零 Neumann 条件使总热量守恒。若使用 \(-\kappa u_x(t,0)+h_0u(t,0)=0\)、\(\kappa u_x(t,\ell)+h_\ell u(t,\ell)=0\)，其中 \(h_0,h_\ell\geq0\)，右端为 \(-h_0u(t,0)-h_\ell u(t,\ell)\)；它对非负温度不大于零。Dirichlet 固定端温却不固定通量，不能由齐次边界推出总热量不变。

\subsection{长时间行为}
\label{b2:subsec:120104}

\Needspace{6\baselineskip}
\begin{proposition}[趋向零模态]\label{b2:prop:12-long-time}
设 \(I=(a,b)\) 为有限区间，\(p,w,q\in C([a,b];\mathbb R)\)、\(p,w>0\)，令 \(A\) 为\cref{b2:def:10-sl-operator}中每端取一个实分离 Dirichlet 或 Robin 条件的算子，其作用为 \(Au=w^{-1}(-(pu')'+qu)\)，\(H=L^2(I,w\,\mathrm dx;\mathbb C)\)，并假设 \(A\geq0\)。对 \(f\in H\)，令 \(u(t)=e^{-tA}f=\sum_ne^{-\lambda_nt}\langle f,e_n\rangle_He_n\)，其中 \((e_n,\lambda_n)\) 是按特征值递增排列的标准正交谱基。\(P_0\) 为到 \(\ker A\) 的正交投影；若没有零特征值则 \(P_0=0\)。记最小正特征值为 \(\lambda_+\)。则对 \(t\geq0\)，
\[
 \|u(t)-P_0f\|_H\leq e^{-\lambda_+t}\|f-P_0f\|_H.
\]
对给定初值 \(f\)，热解趋于零当且仅当 \(P_0f=0\)。所有 \(f\in H\) 的热解都趋于零，当且仅当 \(\ker A=\{0\}\)，这在当前非负正则有限区间条件下等价于 \(\lambda_1>0\)。特别地，对 \(I=(0,\ell)\)、\(A=-\kappa\partial_x^2\)、\(\ell,\kappa>0\)，零 Dirichlet 情形 \(P_0f=0\)，零 Neumann 情形 \(P_0f=\overline f=\ell^{-1}\int_0^\ell f\)；两者 \(\lambda_+=\kappa\pi^2/\ell^2\)。
\end{proposition}
\begin{proof}
将零模态从展开中分离，剩余每个系数的平方都乘以不大于 \(e^{-2\lambda_+t}\) 的因子，求和即得。最小正特征值存在由\cref{b2:thm:10-sl-spectral}保证。于是 \(u(t)\to P_0f\)，给定初值趋零的判据随之成立。若要求所有初值都趋零，取非零 \(f\in\ker A\) 就会得到不动解 \(u(t)=f\)，所以必须核为零；反向由已得估计成立。谱从最低特征值开始且趋于正无穷，非负性下核为零恰好意味着 \(\lambda_1>0\)。纯热 Dirichlet、Neumann 的结论由各自显式谱得到。
\end{proof}

对上述常系数 Dirichlet、Neumann 区间，这还给出 \(t\geq1\) 时的一致收敛：先演化至 \(t-1\)，再由固定时间一的核把 \(L^2\) 误差送到一致范数；Cauchy--Schwarz 与正时间可和尾给出固定的 \(L^2\to L^\infty\) 界。若初值恰与最低正模态正交，衰减速度由下一项非零系数决定。出现负特征值时，相应非零初始系数会产生指数增长，不能继续使用本命题。

\ProseParagraphBreak

对于一般连续系数，判断是否趋于零应检查整个零空间。一个容易核验的充分条件是：\(q\geq0\)、所有 Robin 系数 \(h_s\geq0\)，且至少有一个 Dirichlet 端点或一个严格正的 Robin 系数。事实上，若 \(v\in\ker A\)，则
\[
 0=\langle Av,v\rangle_H
  =\int_a^b\bigl(p|v'|^2+q|v|^2\bigr)\,\mathrm{d}x
    +\sum_{s\in R}h_s|v(s)|^2.
\]
各项非负，\(p>0\) 使 \(v'=0\) 几乎处处，故 \(v\) 为常数。Dirichlet 端点的零值或严格正的 Robin 边界项迫使这个常数为零。因此核为零，由命题得到所有初温的指数衰减。这个论证用到了 \(q\geq0\)，不能仅留下正 Robin 条件而删去对内部系数的控制。

下面的例子说明原因。取 \(I=(0,1)\)、\(p=w=1\)，令
\[
 v_0(x)=1+x(1-x)>0,\qquad
 q(x)=\frac{v_0''(x)}{v_0(x)}=-\frac2{1+x(1-x)},
\]
并在两端取严格正 Robin 系数 \(h_0=h_1=1\)，即
\[
 -v'(0)+v(0)=0,\qquad v'(1)+v(1)=0.
\]
函数 \(v_0\) 满足这两个条件，而且 \(Av_0=-v_0''+qv_0=0\)。还需核验这个例子没有违反非负性假设。对任意 \(v\in H^1(0,1;\mathbb C)\)，令 \(z=v/v_0\)。由于 \(v_0\) 光滑且有正下界，\(z\in H^1(0,1)\)，绝对连续函数的乘积法则给出
\[
 |v'|^2+\frac{v_0''}{v_0}|v|^2
 =v_0^2|z'|^2+\bigl(v_0v_0'|z|^2\bigr)'.
\]
对 \(x\) 积分，导数项成为
\[
 \left[\frac{v_0'}{v_0}|v|^2\right]_0^1
 =-|v(1)|^2-|v(0)|^2,
\]
恰好抵消能量形式中的两个正边界项。因此
\begin{equation}\label{b2:eq:12-robin-zero-mode-factorization}
 \begin{aligned}
 \mathfrak a(v,v)
 &=\int_0^1\bigl(|v'|^2+q|v|^2\bigr)\,\mathrm{d}x
       +|v(0)|^2+|v(1)|^2\\
 &=\int_0^1v_0^2\left|\left(\frac v{v_0}\right)'\right|^2\,\mathrm{d}x
 \geq0.
 \end{aligned}
\end{equation}
所以 \(A\geq0\)，却仍有非恒定零模态。取初温 \(f=v_0\)，就得到永不衰减的解 \(u(t,x)=v_0(x)\)。正 Robin 系数排除了非零常数零模态，但允许 \(q\) 有负值时，不能据此排除所有零模态。

\ProseParagraphBreak

最后回到 \((0,\ell)\) 上 \(u_t=\kappa u_{xx}\)、\(\kappa,\ell>0\) 的纯热模型，把边界的几种作用放在一起比较。周期边界把两个端点耦合为
\[
 u(t,0)=u(t,\ell),\qquad u_x(t,0)=u_x(t,\ell),
\]
因而不属于前面的分离边界条件，不能直接套用其特征值单重的编号。对 \(f\in L^2((0,\ell);\mathbb C)\)，将区间视为周长 \(\ell\) 的圆，独立使用圆周 Fourier 理论：令
\[
 e_k(x)=\ell^{-1/2}e^{2\pi\mathrm{i}kx/\ell},\qquad
 f_k=\int_0^\ell f(x)\overline{e_k(x)}\,\mathrm{d}x,
 \qquad k\in\mathbb Z.
\]
周期热解为
\begin{equation}\label{b2:eq:12-periodic-heat-series}
 u_{\mathrm{per}}(t,x)=\sum_{k\in\mathbb Z}
       f_ke^{-\kappa(2\pi k/\ell)^2t}e_k(x),\qquad t\geq0.
\end{equation}
Parseval 等式及控制收敛使它在零时以 \(L^2\) 范数取得初值；任意正时间带上的指数尾又允许所需的逐项微分。零频率给出 \(\overline f=\ell^{-1}\int_0^\ell f\)，其余频率至少为 \(2\pi/\ell\)，所以
\[
 \|u_{\mathrm{per}}(t)-\overline f\|_2
 \leq e^{-4\kappa\pi^2t/\ell^2}\|f-\overline f\|_2.
\]
周期边界的两个端点通量相消，因此总热量守恒。保正性也可独立看出：周期化 Gaussian
\[
 P_\ell(t,z)=\sum_{m\in\mathbb Z}
  \frac1{\sqrt{4\pi\kappa t}}
  e^{-(z+m\ell)^2/(4\kappa t)},\qquad t>0,
\]
严格为正，且 \(\int_0^\ell P_\ell(t,z)\,\mathrm{d}z=1\)。把 Fourier 系数的积分按周期换元，再用 Gaussian 的变换，得到其第 \(k\) 个系数为 \(\ell^{-1}e^{-\kappa(2\pi k/\ell)^2t}\)；这里系数采用 \(\ell^{-1}\int_0^\ell P_\ell(t,z)e^{-2\pi\mathrm{i}kz/\ell}\,\mathrm{d}z\) 的非单位约定。因此式\eqref{b2:eq:12-periodic-heat-series}也等于 \(\int_0^\ell P_\ell(t,x-y)f(y)\,\mathrm{d}y\)，非负初温产生非负解。

分离边界的纯热模型同样保正，但原因并非仅为谱非负。对实值解令 \(u_-=\max\{-u,0\}\)，在 Dirichlet 端点有 \(u_-=0\)，在 Robin 端点使用 \(-\kappa u_x(0)+h_0u(0)=\kappa u_x(\ell)+h_\ell u(\ell)=0\)，其中 \(h_0,h_\ell\geq0\)，Neumann 对应 \(h_s=0\)。负部属于 \(H^1\)，且其弱导数为 \(-u_x\mathbf1_{\{u<0\}}\)。对正时间强解，将方程与负部配对并分部积分，得到
\[
 \frac12\frac{\mathrm{d}}{\mathrm{d}t}\|u_-\|_2^2
 =-\kappa\|(u_-)_x\|_2^2
      -\sum_{s\in R}h_s|u_-(t,s)|^2\leq0.
\]
时间链式法则可由 \(s\mapsto s_-^2/2\) 的 Lipschitz 导数及 \(L^2\) 时间可微性验证；空间分部积分只用通量和负部的一阶弱导数。从 \(\varepsilon>0\) 积分到 \(t\)，再用 \(u(\varepsilon)\to f\) 及负部映射的 \(L^2\) 连续性令 \(\varepsilon\downarrow0\)，就得到非负初温始终非负。\secref{b2:sec:1203}将对负部方法及最大值原理作进一步说明。

\begin{table}[htbp]
\centering
\small
\setlength{\tabcolsep}{4pt}
\caption[纯热方程的边界、热量与长期平衡]{纯热方程的边界、热量与长期平衡。固定 \(\ell,\kappa>0\)，方程为 \(u_t=\kappa u_{xx}\)，初温 \(f\in L^2((0,\ell);\mathbb R)\)；总热量指 \(\int_0^\ell u\)，长期极限按 \(L^2\) 范数理解，\(\overline f=\ell^{-1}\int_0^\ell f\)。保正指 \(f\geq0\) 时 \(u(t)\geq0\)。Robin 采用正文的含 \(\kappa\) 通量约定。}
\label{b2:tab:12-boundary-heat-equilibrium}
\begin{tabularx}{\linewidth}{@{}>{\raggedright\arraybackslash}p{3.5cm}>{\raggedright\arraybackslash}X>{\raggedright\arraybackslash}p{1.1cm}>{\raggedright\arraybackslash}p{1.8cm}@{}}
\toprule
边界条件 & 总热量 & 保正 & 长期极限 \\
\midrule
两端零 Dirichlet & 一般不守恒 & 是 & \(0\) \\
两端零 Neumann & 守恒 & 是 & \(\overline f\) \\
Robin：\(h_0,h_\ell\geq0\)，至少一端严格正 & 非负初温时不增 & 是 & \(0\) \\
周期边界 & 守恒 & 是 & \(\overline f\) \\
\bottomrule
\end{tabularx}
\end{table}

\Cref{b2:tab:12-boundary-heat-equilibrium}只比较无内部反应的纯热方程。质量守恒、保正性与长期趋零是三个不同的问题：加入 \(q(x)u\) 后，须分别重新检查加权热量平衡、解的符号及 \(\ker A\) 和谱下界，不能仅凭某一个边界系数作判断。

% !TeX root = ../../Book2.tex
\section{Fourier 变换法}
\label{b2:sec:1202}

在全空间中，空间平移不会改变方程，离散特征级数因而换成连续频率积分。变换负责找出核，卷积与逼近恒等负责证明粗初值确实演化成解。外力也可逐时叠加，但积分公式本身并不自动赋予它任意阶的经典导数。

\subsection{热核}
\label{b2:subsec:120201}

全空间没有端点，也没有正的最低允许频率。Fourier 变换把每个连续频率单独演化。先对足够光滑且衰减的解形式计算，按照本卷的 \(2\pi\) 约定有
\[
 \partial_t\widehat u(t,\xi)
 =-4\pi^2\kappa|\xi|^2\widehat u(t,\xi),
 \quad\widehat u(t,\xi)=e^{-4\pi^2\kappa t|\xi|^2}\widehat f(\xi).
\]
\cref{b2:thm:03-gaussian}已经计算了这个乘子的逆变换。沿用第二编的热核记号，在 \(\mathbb R^d\) 上写作
\begin{equation}\label{b2:eq:12-gauss-heat}
 H_t(x)=(4\pi t)^{-d/2}e^{-|x|^2/(4t)},
 \quad t>0,
 \quad\widehat H_t(\xi)=e^{-4\pi^2t|\xi|^2}.
\end{equation}
方程扩散系数为 \(\kappa\) 时使用 \(H_{\kappa t}\)，不另换 Fourier 归一化。这个核非负、总积分为一；其空间宽度约为 \(\sqrt t\)，高度约为 \(t^{-d/2}\)。由直接微分有 \(\partial_tH_{\kappa t}=\kappa\Delta H_{\kappa t}\)。

\ProseParagraphBreak

这些计算可以独立验证卷积候选，不要求先假定未知解的变换存在。对于粗初值，求解的逻辑次序是先用核定义 \(u\)，再证明可微性和初值收敛；前面的频率微分只负责找到候选。第二编已经证明 Gauss 核的逼近恒等性质，下面把其参数取为 \(\sqrt{4\pi\kappa t}\)。

\subsection{卷积表示}
\label{b2:subsec:120202}

\Needspace{6\baselineskip}
\begin{theorem}[全空间热核解]\label{b2:thm:12-heat-convolution}
设整数 \(d\geq1\)、\(\kappa>0\)，\(H_s(x)=(4\pi s)^{-d/2}e^{-|x|^2/(4s)}\)（\(s>0\)、\(x\in\mathbb R^d\)）。设 \(1\leq p\leq\infty\)、\(f\in L^p(\mathbb R^d;\mathbb C)\)。则
\begin{equation}\label{b2:eq:12-heat-convolution}
 u(t,x)=\int_{\mathbb R^d}H_{\kappa t}(x-y)\cdot f(y)\,\mathrm{d}y
\end{equation}
对 \(t>0\)、\(x\in\mathbb R^d\) 有意义，且在 \((0,\infty)\times\mathbb R^d\) 上光滑，满足 \(u_t=\kappa\Delta u\)。当 \(t\downarrow0\) 且 \(p<\infty\) 时，\(u(t)\to f\) 于 \(L^p\)。若选定的代表 \(\widetilde f\) 有界一致连续，则 \(u(t)\to\widetilde f\) 在一致范数中成立；若选定 \(f\) 的一个有界代表 \(\widetilde f\)，且它在 \(x_0\in\mathbb R^d\) 连续，则 \(u(t,x_0)\to\widetilde f(x_0)\)。
\end{theorem}
\begin{proof}
热核及其每阶时空导数属于每个 \(L^r\)、\(1\leq r\leq\infty\)。固定 \(t\) 的正紧区间与 \(x\) 的紧集后，每个导数可由一个平移后仍快速衰减的函数统一控制。对 \(f\in L^p\)，Hölder 不等式给出绝对可积性；对差商使用同一控制，可以把每一阶导数移入积分。因此卷积光滑，且核的微分等式传给 \(u\)。这里在 \(p=1\) 用核导数的一致界，在 \(p=\infty\) 用其绝对可积性，两个端点都被覆盖。

有限 \(p\) 与一致连续代表的情形由\cref{b2:thm:03-approximate-identity}得到。卷积不因初值在零测集上的修改而改变。对选定代表的连续点，把 \(\widetilde f(x_0-y)-\widetilde f(x_0)\) 的积分分成 \(|y|<\delta\) 与补集。近处由连续性控制，远处由 \(2\sup_x|\widetilde f(x)|\) 乘热核尾质量控制，先选 \(\delta\) 再令 \(t\downarrow0\) 即得。这样取得的点值是所选连续代表的点值，而不是给 \(L^p\) 等价类额外指定一个任意值。
\end{proof}

初值也可以是一点质量，此时解是 \(H_{\kappa t}(x-y_0)\)，在 \(t\downarrow0\) 时只以测度或分布意义趋于 \(\delta_{y_0}\)。为了区别空间核与时空基本解，令 \(E(t,x)=H_{\kappa t}(x)\) 当 \(t>0\)，在 \(t<0\) 令 \(E=0\)。那么
\[
 (\partial_t-\kappa\Delta_x)E=\delta_{(0,0)}
\]
在时空分布中成立。确实，对紧支集光滑 \(\varphi(t,x)\)，从 \(t=\varepsilon\) 起分部积分，热方程抵消内部项，留下 \(\int H_{\kappa\varepsilon}(x)\cdot\varphi(\varepsilon,x)\,\mathrm{d}x\)。它由逼近恒等趋于 \(\varphi(0,0)\)。同时 \(E\) 在有限时间带的空间积分为该带长度，所以零时附近的局部可积性也已得到。

\ProseParagraphBreak

对非齐次方程，设整数 \(d\geq1\)、\(\kappa,T>0\)、\(1\leq p<\infty\)，\(f\in L^p(\mathbb R^d;\mathbb C)\)、\(F\in C([0,T];L^p(\mathbb R^d;\mathbb C))\)。由逐时叠加得到
\begin{equation}\label{b2:eq:12-duhamel}
 u(t)=H_{\kappa t}*f+
 \int_0^tH_{\kappa(t-s)}*F(s)\,\mathrm{d}s.
\end{equation}
在积分的右端按 \(H_0*F(t)=F(t)\) 理解。被积函数是连续的 \(L^p\) 值函数，积分可作为 Riemann 和的范数极限；Young 不等式给出
\[
 \left\|\int_0^tH_{\kappa(t-s)}*F(s)\,\mathrm{d}s\right\|_p
 \leq\int_0^t\|F(s)\|_p\,\mathrm{d}s.
\]
这保证源项趋于零初值。先对光滑紧支集的时空源项求导，再以连续 \(L^p\) 范数逼近一般 \(F\)，便在分布中得到 \(u_t-\kappa\Delta u=F\)。逼近可将时间分成有限小段，对有限个空间取值分别磨光，再以时间连续插值连接；上式使解同时收敛，测试函数的积分配对也连续。对于这里的一般连续 \(L^p\) 值源项，方程按分布意义成立；经典二阶空间正则性还需进一步的条件。

\subsection{高维热方程}
\label{b2:subsec:120203}

高维热核是各坐标一维核的乘积，但高维扩散的幅度不能从一维直观中省略维数。换元 \(x=\sqrt t\,z\) 给出
\[
 \|H_t\|_r=C_{d,r}t^{-\frac d2(1-1/r)},
 \quad1\leq r\leq\infty.
\]
例如有限质量初温满足 \(\|u(t)\|_\infty\leq(4\pi\kappa t)^{-d/2}\|f\|_1\)。没有任何质量损失也能使温度峰值降低，因为同一热量散布到越来越大的区域。

\ProseParagraphBreak

方向上的扩散速率也可以不同。若整数 \(d\geq1\)，\(B\) 是实对称正定的 \(d\times d\) 常矩阵，方程为 \(u_t=\sum_{i,j=1}^dB_{ij}\partial_i\partial_ju\)，则 \(t>0\) 时的相应核是
\[
 H_t^B(x)=\dfrac1{(4\pi t)^{d/2}\sqrt{\det B}}
 \exp\left(-\dfrac{x^{\mathsf T}B^{-1}x}{4t}\right).
\]
正交对角化 \(B\) 后，各方向分别使用一维扩散核，换元的行列式给出归一化因子。故核仍非负、质量为一，变换为 \(e^{-4\pi^2t\xi^{\mathsf T}B\xi}\)。正定性使每个方向都发生扩散。若 \(B\) 仅半正定且有零特征值，\(B^{-1}\) 与 \((\det B)^{-1/2}\) 就没有上述意义，不能在函数密度公式中直接把某个特征值置零。

\ProseParagraphBreak

例如在 \(\mathbb R^2\) 上取 \(B=\operatorname{diag}(1,0)\)，方程成为 \(u_t=u_{xx}\)，没有 \(y\) 方向的扩散。此时传播核是乘积测度
\[
 H_t^{(1)}(x)\otimes\delta_0(y),\qquad
 H_t^{(1)}(x)=(4\pi t)^{-1/2}e^{-x^2/(4t)},\quad t>0,
\]
它支撑在直线 \(y=0\) 上，不是相对于二维 Lebesgue 测度的可积函数密度。对 \(f\in L^2(\mathbb R^2;\mathbb C)\)，传播只沿 \(x\) 作卷积：
\[
 u(t,x,y)=\int_{\mathbb R}H_t^{(1)}(x-z)f(z,y)\,\mathrm{d}z.
\]
Fubini 定理使几乎每个 \(y\) 对应的初值切片属于 \(L^2(\mathbb R_x)\)，故这一维卷积有意义，并在分布意义下满足 \(u_t=u_{xx}\)。但它不提高横向正则性：取非负非零 \(\phi\in C_c^\infty(\mathbb R;\mathbb R)\)，令 \(f(x,y)=\phi(x)\mathbf1_{(0,1)}(y)\)，就有
\[
 u(t,x,y)=\bigl(H_t^{(1)}*\phi\bigr)(x)\mathbf1_{(0,1)}(y).
\]
前一个因子在 \(t>0\) 时光滑且处处为正，\(y=0,1\) 处的跳跃却仍保留。一般半正定矩阵经正交对角化后也如此：正特征值方向使用 Gaussian 核，零特征值方向保留点质量，不能从正方向的磨光推出全空间光滑化。

\ProseParagraphBreak

这一正定二次型公式可与 John K. Hunter、Bruno Nachtergaele 的《Applied Analysis》\cite[pp. 303--304]{Hunter2001Applied}对照，其变换采用不同的 \(2\pi\) 归一化。本卷已经在\secref{b2:sec:0502}换算过常数，这里直接使用一致的频率约定。

% !TeX root = ../../Book2.tex
\section{能量泛函}
\label{b2:sec:1403}

同一个偏微分方程可以通过能量的驻点条件出现，但“能量”一词并不总表示待极小化的非负量。椭圆方程常对应空间能量的极小值；波动方程对应时空作用量的驻点，同时另有随时间守恒的正能量。本节用这两个模型把几种不同作用区分开。

\subsection{Dirichlet 能量}
\label{b2:subsec:140301}

设 \(\Omega\subset\mathbb R^d\) 是有界开集，记 \(H_0^1(\Omega)=W_0^{1,2}(\Omega)\)。第一卷定义 22.2.1 将它定义为 \(C_c^\infty(\Omega)\) 在 \(H^1\) 中的闭包，这一定义不要求边界光滑。由第一卷定理 24.1.2 的零边界 Poincaré 不等式，\(\lVert v\rVert_2\le C_P\lVert\nabla v\rVert_2\)，所以 \(\lVert\nabla v\rVert_2\) 是这个空间上的等价 Hilbert 范数。

\ProseParagraphBreak

设 \(\Omega\subset\mathbb R^d\) 是有界开集，给定实值 \(f\in L^2(\Omega)\) 和实值 \(g\in H^1(\Omega)\)，把允许函数写为 \(g+H_0^1(\Omega)\)，考虑 \term[dirichlet-nengliang]{Dirichlet 能量}[Dirichlet energy]及线性外力项组成的泛函
\[
 E(u)=\frac12\int_\Omega|\nabla u|^2\,\mathrm dx-\int_\Omega fu\,\mathrm dx.
\]
仿射空间的写法精确表达“\(u\) 与 \(g\) 具有相同边界数据”：即使边界粗糙到尚未定义逐点边界迹，也可以使用它。

\Needspace{6\baselineskip}
\begin{proposition}[二次能量的驻点与极小点]
\label{b2:prop:14-dirichlet-minimum}
设整数 \(d\geq1\)，\(\Omega\subset\mathbb R^d\) 为非空有界开集，\(f\in L^2(\Omega;\mathbb R)\)、\(g\in H^1(\Omega;\mathbb R)\)，\(E(v)=\frac12\int_\Omega|\nabla v|^2-\int_\Omega fv\)。\(u\in g+H_0^1(\Omega;\mathbb R)\) 是 \(E\) 在该仿射空间上的极小点，当且仅当
\[
 \int_\Omega\nabla u\cdot\nabla h\,\mathrm dx=\int_\Omega fh\,\mathrm dx
 \qquad(h\in H_0^1(\Omega)).
\]
这样的极小点至多有一个，且对每个 \(h\in H_0^1(\Omega)\) 都有
\[E(u+h)-E(u)=\tfrac12\lVert\nabla h\rVert_2^2.\]
\end{proposition}
\begin{proof}
Cauchy--Schwarz 不等式保证所有积分有限。直接展开平方得
\[
 E(u+h)-E(u)=\int_\Omega\nabla u\cdot\nabla h-\int_\Omega fh
              +\tfrac12\lVert\nabla h\rVert_2^2.
\]
其中前两项构成连续线性泛函，最后一项是二阶余项，故一阶导数正是前两项。极小点的双侧变分为零，得到所列积分等式。反之，积分等式使能量差等于非负的梯度平方，所以 \(u\) 为全局极小点。若另一点 \(u+h\) 也极小，则 \(\nabla h=0\)；Poincaré 不等式推出 \(h=0\)。
\end{proof}

这一命题暂时没有断言极小点存在；它说明一旦构造出驻点，就无须再排查鞍点。二次展开还给出误差的精确解释：任何满足同一边界数据的近似函数，其能量超额恰为梯度误差平方的一半。与一般的第二变分局部判别相比，这是全局恒等式，原因在于能量严格二次且对允许方向一致凸。

\ProseParagraphBreak

在区间 \((0,1)\) 上取 \(f=1\)、\(g=0\)。函数 \(u(x)=x(1-x)/2\) 满足 \(-u''=1\) 和零端点条件；对光滑紧支集扰动分部积分，再用密度推广到 \(H_0^1\)，即可验证积分等式。因此它是唯一极小点，能量为 \(-1/24\)。负能量并不矛盾：梯度项非负，外力项却可以降低总泛函值；强制性要求的是大范数下的增长，而不是每一点的非负性。

\subsection{波动作用量与守恒能量}
\label{b2:subsec:140302}

设整数 \(d\geq1\)，\(\Omega\subset\mathbb R^d\) 为非空有界 \(C^1\) 开集，\(c,T>0\)，并考虑在 \([0,T]\times\overline\Omega\) 的开邻域内为 \(C^2\)、在侧边界取零值的实函数 \(u(t,x)\)。时空作用量为
\[
 \mathcal A(u)=\frac12\int_0^T\int_\Omega
       \bigl(|\partial_tu|^2-c^2|\nabla_xu|^2\bigr)\,\mathrm dx\,\mathrm dt.
\]
对在初、末时刻和空间边界均为零的扰动 \(h\)，时、空分别分部积分，得到
\[
 D\mathcal A(u)h=\int_0^T\int_\Omega
       (-\partial_t^2u+c^2\Delta_xu)h\,\mathrm dx\,\mathrm dt.
\]
所以驻点满足波动方程 \(\partial_t^2u-c^2\Delta u=0\)。这里固定初末值是推导变分方程时对扰动的限制；实际波动初值问题仍由初位移和初速度确定，不能把两种问题混为一谈。

\ProseParagraphBreak

作用量的二次型有正负两部分，一般并不在解处取得极小值。例如取乘积扰动 \(h(t,x)=a(t)b(x)\)，其中两因子均光滑紧支集，则
\[
 2\mathcal A(h)=\lVert a'\rVert_2^2\lVert b\rVert_2^2
      -c^2\lVert a\rVert_2^2\lVert\nabla b\rVert_2^2.
\]
固定非零 \(b\) 后，取 \(a_N(t)=\eta(t)\sin(Nt)\)，其中 \(0\ne\eta\in C_c^\infty(0,T)\)。展开正弦平方并用\cref{b2:thm:03-riemann-lebesgue}，有 \(\|a_N\|_2^2\to\|\eta\|_2^2/2\) 及 \(N^{-2}\|a_N'\|_2^2\to\|\eta\|_2^2/2\)，所以作用量为正。反之，固定非零 \(a\)，取 \(b_N(x)=\psi(x)\sin(Nx_1)\)，其中 \(0\ne\psi\in C_c^\infty(\Omega)\)。相同的平方展开给出 \(\|b_N\|_2^2\to\|\psi\|_2^2/2\) 及 \(N^{-2}\|\nabla b_N\|_2^2\to\|\psi\|_2^2/2\)，所以作用量为负。若 \(u\) 为驻点，精确二次展开给出 \(\mathcal A(u+th)-\mathcal A(u)=t^2\mathcal A(h)\)。故可把两类扰动都缩得任意小，在每个驻点附近仍有增加和降低作用量的方向。

\Needspace{6\baselineskip}
\begin{proposition}[波动能量守恒]
\label{b2:prop:14-wave-energy}
设整数 \(d\geq1\)，\(c,T>0\)，\(\Omega\subset\mathbb R^d\) 为非空有界 \(C^1\) 开集，实函数 \(u\) 在 \([0,T]\times\overline\Omega\) 的某开邻域内属于 \(C^2\)，满足 \(u_{tt}-c^2\Delta u=0\) 于 \((0,T)\times\Omega\)，且 \(u(t,y)=0\) 对 \(0\leq t\leq T\)、\(y\in\partial\Omega\) 成立。则
\[
 \mathcal E(t)=\frac12\int_\Omega
          \bigl(|u_t(t,x)|^2+c^2|\nabla u(t,x)|^2\bigr)\,\mathrm dx
\]
与时间无关。
\end{proposition}
\begin{proof}
对时间求导并对空间分部积分，得
\[
 \mathcal E'(t)=\int_\Omega u_t(u_{tt}-c^2\Delta u)\,\mathrm dx
           +c^2\int_{\partial\Omega}u_t\partial_\nu u\,\mathrm dS=0.
\]
内部项由方程为零，边界项由对恒定零边界条件求时间导数所得 \(u_t=0\) 而消失。
\end{proof}

守恒能量中空间项带正号，作用量中带负号。前者衡量一个时刻的状态，后者比较整段时空运动。对同样的区域、实值 \(C^2\) 解及恒定零边界，若方程改为 \(u_{tt}-c^2\Delta u+\gamma u_t=f\)，其中 \(\gamma\geq0\)、\(f\) 连续，则同一计算给出
\[
 \mathcal E'(t)=\int_\Omega f(t,x)u_t(t,x)\,\mathrm{d}x
                    -\gamma\|u_t(t)\|_2^2.
\]
这些等式将外界做功和耗散分开记录。\cref{b2:tab:14-action-energy}集中比较作用量与能量；表中的守恒结论始终包含零边界且无外力、无阻尼的条件。

\begin{table}[htbp]
\centering
\caption[波动作用量与守恒能量的区别]{波动作用量与守恒能量的区别。取整数 \(d\geq1\)、\(c,T>0\)、非空有界 \(C^1\) 开集 \(\Omega\subset\mathbb R^d\)，函数实值且光滑到足以完成正文的微分与分部积分，侧边界为恒定零值。作用量列的变分还固定初、末时刻的函数值；能量列讨论满足波动方程的实际演化。}
\label{b2:tab:14-action-energy}
\small
\begin{tabularx}{\textwidth}{@{}p{5.2em}XX@{}}
\toprule
比较项目 & 时空作用量 \(\mathcal A\) & 守恒能量 \(\mathcal E(t)\)\\
\midrule
积分范围 & 整段 \((0,T)\times\Omega\) & 固定时刻的 \(\Omega\)\\
空间梯度项 & \(-c^2|\nabla u|^2/2\) & \(+c^2|\nabla u|^2/2\)\\
数学作用 & 对允许扰动驻定，推出运动方程 & 控制速度与梯度，进而估计解的差\\
极值或演化 & 驻点一般为鞍点 & 无外力、无阻尼时随时间守恒\\
\bottomrule
\end{tabularx}
\end{table}

\subsection{PDE 与极值问题}
\label{b2:subsec:140303}

从微分方程寻找泛函时，最容易检验的是对称性。设 \(H\) 是实 Hilbert 空间，\(a:H\times H\to\mathbb R\) 是有界双线性型，\(\ell\in H^*\)。令
\[J(u)=\tfrac12a(u,u)-\ell(u).\]
直接展开可得
\[DJ(u)h=\tfrac12\bigl(a(u,h)+a(h,u)\bigr)-\ell(h).\]
因此仅当 \(a\) 对称时，驻点条件才是 \(a(u,h)=\ell(h)\)。反对称部分在 \(a(u,u)\) 中完全消失，却可能实质影响方程的解。

\ProseParagraphBreak

例如 \(H=\mathbb R^2\)，令 \(Au=(u_1-u_2,u_1+u_2)\)，\(a(u,h)=(Au)\cdot h\)。尽管 \(a(u,u)=|u|^2\) 严格为正，极小化 \(\tfrac12|u|^2-f\cdot u\) 得到的是 \(u=f\)，而原方程是 \(Au=f\)。两者一般不同。下一章的 Lax--Milgram 定理直接求解有界强制双线性型，无须先把它变成能量；这正是从变分极小化过渡到一般弱解方法时必须增加的能力。

\ProseParagraphBreak

另一个区别是驻点与动力学。若 \(E(u)=\tfrac12\int|\nabla u|^2\)，形式上的 \(L^2\) 梯度为 \(-\Delta u\)。沿最陡下降运动 \(u_t=-\operatorname{grad}E(u)\)，得到热方程 \(u_t=\Delta u\)，而平衡点满足 Laplace 方程。对光滑零边界解，分部积分给出
\[\frac{\mathrm d}{\mathrm dt}E(u(t))=-\int_\Omega|u_t|^2\,\mathrm dx\le0.\]
这项恒等式说明热演化趋向降低能量，却不在这里自动证明长时间收敛；还需要轨道紧性、极限点的识别等信息。变分结构可以提供估计和候选极限，但不能代替这些额外步骤。

% !TeX root = ../../Book2.tex
\section{直接方法}
\label{b2:sec:1404}

先推导 Euler--Lagrange 方程再求解，有时会把存在性问题变得更难。直接方法反过来考察能量本身：取能量趋近下确界的序列，控制它的大小，抽取极限，再证明极限没有丢失边界条件和能量下界。第一卷定理 9.5.9 已在自反 Banach 空间中证明这一抽象原则。本节不重复该证明，而逐项核验积分能量满足条件的具体原因。

\subsection{强制性}
\label{b2:subsec:140401}

在 Banach 空间 \(X\) 的约束集 \(K\) 上，泛函的强制性指
\[J(v)\longrightarrow+\infty\qquad(v\in K,\ \lVert v\rVert_X\to\infty).\]
它保证能量有界的候选者不能逃到范数无穷远处。在零边界空间 \(H_0^1(\Omega)\) 上，Poincaré 不等式和 Cauchy--Schwarz 不等式给出
\[
 \frac12\lVert\nabla v\rVert_2^2-\int_\Omega fv
 \ge\frac12\lVert\nabla v\rVert_2^2-C_P\lVert f\rVert_2\lVert\nabla v\rVert_2.
\]
右端是开口向上的二次式，故这项能量强制。若允许非零边界数据，写成 \(u=g+v\)，交叉项 \(\int\nabla g\cdot\nabla v\) 仍只按 \(\lVert\nabla v\rVert_2\) 线性增长，因此同样得到强制性。

\ProseParagraphBreak

边界限制不能从这个论证中删除。在 \(H^1(\Omega)\) 全空间中，常函数的梯度能量为零，故单靠梯度无法控制范数。若 \(\int_\Omega f\ne0\)，沿适当符号的常函数甚至能让总能量趋向 \(-\infty\)。Neumann 问题因此需要相容条件和一个固定常数的归一化；零均值空间上的 Poincaré 不等式才是相应替代物。

\subsection{极小化序列}
\label{b2:subsec:140402}

若 \(m=\inf_KJ\) 是有限实数，可以按下确界的定义选取 \(u_n\in K\)，满足 \(m\le J(u_n)<m+1/n\)。这样的序列称为极小化序列。选取过程既不要求驻点已存在，也不要求每个 \(u_n\) 近似满足 Euler--Lagrange 方程。

\ProseParagraphBreak

控制能量后，强制性使序列有界，但有界不等于强收敛。无限维 Hilbert 空间中的标准正交序列有界，却没有强收敛子列；它仅弱收敛到零。这就是直接方法引入弱拓扑的实质理由。另一方面，若约束集不是弱闭的，弱极限可能不再可行；若能量不是弱下半连续的，极限能量可能突然升高。序列有界、约束保留与能量比较分别承担不同任务，不能用一句“由紧性”全部略去。

\ProseParagraphBreak

二次能量还有一个可直接看见的特殊现象。若 \(u_n\) 和 \(u_k\) 满足同一边界数据，则
\[
 E\left(\frac{u_n+u_k}{2}\right)
 =\frac{E(u_n)+E(u_k)}2-\frac18\lVert\nabla(u_n-u_k)\rVert_2^2.
\]
左端不小于下确界 \(m\)。因此极小化序列事实上在梯度范数下为 Cauchy 序列，借完备性可以直接得到强极限。一般非二次能量不一定有这个精确恒等式，弱紧性的方法因而更具普遍性。

\subsection{弱紧性}
\label{b2:subsec:140403}

设 \(1<p<\infty\)。空间 \(W^{1,p}(\Omega)\) 可以等距嵌入 \((L^p(\Omega))^{d+1}\)，其像由弱导数关系确定，是闭子空间。有限个自反空间的乘积及其闭子空间仍自反；故第一卷定理 9.4.7 给出：每个有界 \(W^{1,p}\) 序列有弱收敛子列。具体而言，沿同一子列有
\[u_n\rightharpoonup u\ \text{于 }L^p,\qquad \partial_i u_n\rightharpoonup\partial_i u\ \text{于 }L^p.\]
弱导数关系可以直接验证：对测试函数 \(\varphi\)，在 \(\int u_n\partial_i\varphi=-\int(\partial_i u_n)\varphi\) 中取极限即可。

\ProseParagraphBreak

仿射边界类 \(g+W_0^{1,p}(\Omega)\) 是强闭凸集，因而弱闭。这里用的是第一卷定理 9.5.2 的闭凸集结论，完全不要求逐点读取边界值。若约束还包含 \(u\ge\psi\) 几乎处处，它也是强闭凸约束：强收敛子列可进一步逐点几乎处处收敛，从而保留不等式；闭凸性再给出弱闭性。

\ProseParagraphBreak

端点 \(p=1\) 与 \(p=\infty\) 不在上述自反性结论内。例如在 \((0,1)\) 上，\(u_n(x)=\min\{nx,1\}\) 的 \(W^{1,1}\) 范数有界，但导数集中到越来越短的边界层中，不能在 \(L^1\) 内取得弱收敛子列。其极限在内部是常数 \(1\)，左端数据却已丢失。处理线性增长能量时往往需要有界变差函数和测度导数；不能把 \(p>1\) 的论证直接搬到这个端点。

\subsection{弱下半连续性}
\label{b2:subsec:140404}

第一卷定理 9.5.8 说明：凸泛函的范数下半连续性等价于弱下半连续性。对积分泛函，一个常用核验顺序是先证强收敛下的下半连续，再证凸性。设 \(F(x,z,p)\ge0\) 对 \((z,p)\) 连续且凸，对 \(x\) 可测，则
\[I(u)=\int_\Omega F(x,u,\nabla u)\,\mathrm dx\]
是允许取 \(+\infty\) 的凸泛函。若 \(u_n\to u\) 于 \(W^{1,p}\)，先抽取实现能量下极限的子列，再从中抽取函数与全部一阶导数几乎处处收敛的子列。连续性与 Fatou 引理给出 \(I(u)\le\liminf I(u_n)\)。因此 \(I\) 强下半连续，继而弱下半连续。

\ProseParagraphBreak

非负性可放宽为统一的可积下界，因为减去该下界不改变论证。又若 \(\ell\in(W^{1,p})^*\)，线性项在弱收敛下连续，所以 \(I-\ell\) 仍弱下半连续。这个证明覆盖了 \(\tfrac1p\int|\nabla u|^p\) 及许多非负凸势能，却没有覆盖任意非凸密度。特别是向量值问题中，关于梯度的普通凸性虽为充分条件，却远非最一般的下半连续性条件；本章不把更深的条件隐含在这里。

\ProseParagraphBreak

以 \(\ell^2\) 上的 \(J(x)=-\lVert x\rVert^2\) 为例，\(J\) 范数连续，但对标准正交序列 \(e_n\rightharpoonup0\) 有 \(J(0)=0>-1=\lim J(e_n)\)。所以范数连续本身不能代替弱下半连续；前面的凸性在改变拓扑时起了决定作用。

\subsection{极小元的存在性}
\label{b2:subsec:140405}

下面把这些核验汇集到一个非线性模型。它既包含 Dirichlet 能量，也说明直接方法并不依赖显式解或线性叠加。

\begin{theorem}[梯度幂能量的极小元]
\label{b2:thm:14-p-energy-minimum}
设 \(\Omega\subset\mathbb R^d\) 是有界开集，\(1<p<\infty\)，\(g\in W^{1,p}(\Omega)\)，\(f\in L^{p'}(\Omega)\)，其中 \(p'=p/(p-1)\)。则
\[
 J(u)=\frac1p\int_\Omega|\nabla u|^p\,\mathrm dx-\int_\Omega fu\,\mathrm dx
\]
在 \(g+W_0^{1,p}(\Omega)\) 上有唯一极小元 \(u\)。它满足
\[
 \int_\Omega|\nabla u|^{p-2}\nabla u\cdot\nabla h\,\mathrm dx
     =\int_\Omega fh\,\mathrm dx\qquad(h\in W_0^{1,p}(\Omega)),
\]
其中在 \(\nabla u=0\) 处约定 \(|\nabla u|^{p-2}\nabla u=0\)。
\end{theorem}
\begin{proof}
令 \(u=g+v\)。由三角不等式和 Poincaré 不等式，\(\lVert\nabla v\rVert_p\le\lVert\nabla u\rVert_p+\lVert\nabla g\rVert_p\)，且
\[\left|\int_\Omega fu\right|\le\lVert f\rVert_{p'}\bigl(\lVert g\rVert_p+C_P\lVert\nabla v\rVert_p\bigr).\]
因此 \(J(g+v)\) 的下界由一个正的 \(\lVert\nabla u\rVert_p^p\) 项减去线性项构成。当 \(\lVert v\rVert_{W^{1,p}}\to\infty\) 时，\(\lVert\nabla u\rVert_p\to\infty\)，而 \(p>1\)，故泛函强制且有有限下界。可行集非空、弱闭；上一小节已核验梯度幂积分弱下半连续，而 \(f\) 的配对弱连续。于是自反性及第一卷定理 9.5.9 给出极小元。

若 \(u_1,u_2\) 都极小，由 \(\xi\mapsto|\xi|^p\) 的严格凸性，除非 \(\nabla u_1=\nabla u_2\) 几乎处处，中点能量就严格更小。故梯度相等；差属于 \(W_0^{1,p}\)，Poincaré 不等式推出差为零。

最后固定 \(h\in W_0^{1,p}\)。沿线段求导时，\(|t|\le1\) 的导数被
\[C_p\bigl(|\nabla u|^{p-1}+|\nabla h|^{p-1}\bigr)|\nabla h|\]
支配；该函数由 Hölder 不等式可积。支配收敛允许交换求导与积分，得第一变分公式。极小点的正负方向都可行，故第一变分为零，即所列等式。
\end{proof}

当 \(p=2\) 时，这就是 Poisson 方程的弱形式。一般 \(p\) 则对应 \(-\operatorname{div}(|\nabla u|^{p-2}\nabla u)=f\)，通常称为 \(p\)-Laplace 方程。上述存在性并没有证明极小元是 \(C^2\) 函数；能写出弱方程与能把它逐点展开，是两个不同层次的结论。

\ProseParagraphBreak

如果再加上非负连续凸势能 \(\int_\Omega\Psi(u)\)，只要某个可行函数具有有限能量，同样的直接方法仍给出极小元；严格凸的梯度项保留唯一性。但要写出 \(\Psi'(u)\) 的弱方程，还须补充可微性和增长条件，确保扰动可积并允许积分下求导。存在性所需的下半连续性通常比推导经典微分方程所需的光滑性弱得多，这也是直接方法的重要优势。

% 本章小结（不计入编号）
\begin{chaptersummary}
第一变分把极小问题转成对全部允许方向成立的积分等式；分部积分把该等式拆成内部 Euler--Lagrange 方程和自然边界条件。允许的方向由约束决定，故固定边界、自由边界及单侧约束产生不同形式。第二变分非负只是必要条件，一致正下界才给出局部极小性的稳定判据。对凸二次能量，精确能量差进一步把驻点与全局唯一极小点联系起来。

\ProseParagraphBreak

直接方法不先求解微分方程，而从极小化序列出发。强制性控制序列，Sobolev 空间的自反性提供弱极限，闭凸约束保留边界数据，弱下半连续性完成能量比较。所获极小元首先属于弱函数空间；它是否更光滑，需要另外的方程估计。下一章将以弱形式直接求解一般双线性问题，尤其处理不能来自对称能量的算子。
\end{chaptersummary}


\BookExercises

下面各题均为本书自拟。题目通过改变空间、边界或能量符号，检验正文结论的适用范围。

% 自拟：无限维正定第二变分不足以保证极小。
\begin{exercise}\label{b2:ex:14-positive-hessian}
在实空间 \(\ell^2\) 上令 \(J(x)=\sum_{n\ge1}x_n^2/n-\lVert x\rVert^4\)。证明 \(J\) 二次连续 Fréchet 可微，零点为驻点且 \(D^2J(0)[h,h]>0\) 对每个 \(h\ne0\) 成立，但零点不是局部极小点。指出第二变分判别的哪项条件不成立。
\end{exercise}
\begin{solution}
有界自伴算子 \(Ax=(x_n/n)\) 给出 \(J(x)=\langle Ax,x\rangle-\langle x,x\rangle^2\)，所以它是连续多项式，特别是二次连续可微。\(DJ(0)=0\)，且 \(D^2J(0)[h,h]=2\sum h_n^2/n>0\)。令 \(x^{(n)}=\sqrt{2/n}\,e_n\)，则 \(\lVert x^{(n)}\rVert\to0\)，而 \(J(x^{(n)})=2/n^2-4/n^2=-2/n^2<0=J(0)\)。在单位向量 \(e_n\) 上第二变分等于 \(2/n\to0\)，故不存在统一正下界 \(c\lVert h\rVert^2\)。
\end{solution}

% 自拟：带权端点问题，应用驻点和严格凸性。
\begin{exercise}\label{b2:ex:14-weighted-string}
设 \(a\in C^1([0,1])\)、\(a\ge a_0>0\)。求 \(\tfrac12\int_0^1a(x)|u'(x)|^2\,\mathrm dx\) 在 \(u(0)=0,u(1)=1\) 的 \(H^1\) 函数中的极小元和极小值，并用能量差证明唯一性。
\end{exercise}
\begin{solution}
记 \(L=\int_0^1a(s)^{-1}\,\mathrm ds\)，令 \(u(x)=L^{-1}\int_0^x a(s)^{-1}\,\mathrm ds\)，则 \(au'=L^{-1}\)。任一零端点扰动 \(h\in H_0^1(0,1)\) 满足 \(\int au'h'=L^{-1}(h(1)-h(0))=0\)，此等式也可由光滑逼近得到。因此能量差为 \(\tfrac12\int a|h'|^2\ge0\)，且等号迫使 \(h'=0\)、\(h=0\)。极小值是 \((2L^2)^{-1}\int a^{-1}=1/(2L)\)。
\end{solution}

% 自拟：边界弹簧与外法向符号。
\begin{exercise}\label{b2:ex:14-boundary-spring}
给定 \(\beta>0\)、\(q\in\mathbb R\)，在 \(u\in H^1(0,1)\)、\(u(0)=0\) 中极小化
\[J(u)=\tfrac12\int_0^1|u'|^2\,\mathrm dx+\tfrac\beta2u(1)^2-qu(1).\]
求极小元，并核验自然边界条件与全局极小性。
\end{exercise}
\begin{solution}
第一变分为 \(\int u'h'+(\beta u(1)-q)h(1)\)，其中 \(h(0)=0\)。内部方程是 \(u''=0\)，右端条件是 \(u'(1)+\beta u(1)=q\)，所以 \(u(x)=qx/(1+\beta)\)。对任意允许扰动，交叉项消失，能量差为 \(\tfrac12\int|h'|^2+\tfrac\beta2h(1)^2\)。它非负，等号时 \(h=0\)，故所得解是唯一全局极小元。
\end{solution}

% 自拟：积分约束产生常数乘子，含显式解。
\begin{exercise}\label{b2:ex:14-mass-constraint}
给定 \(M\in\mathbb R\)，在 \(H_0^1(0,1)\) 中，附加 \(\int_0^1u=M\)，极小化 \(\tfrac12\int_0^1|u'|^2\)。不用抽象乘子定理，求出唯一极小元和极小值。
\end{exercise}
\begin{solution}
取 \(u(x)=6Mx(1-x)\)，则端点为零且积分为 \(M\)。允许扰动 \(h\) 满足零端点和零积分；分部积分给出 \(\int u'h'=-\int u''h=12M\int h=0\)。能量差为 \(\tfrac12\lVert h'\rVert_2^2\)，故它是唯一极小元。又 \(\int_0^1(1-2x)^2\,\mathrm dx=1/3\)，极小值为 \(6M^2\)。常数 \(-u''=12M\) 正是积分约束所允许留下的乘子；若错误地令所有紧支集方向都可行，就会误得 \(u''=0\)。
\end{solution}

% 自拟：临界负势能与最优Poincare常数。
\begin{exercise}\label{b2:ex:14-negative-potential}
在 \(H_0^1(0,\pi)\) 上考察 \(J_\lambda(u)=\tfrac12\int_0^\pi(|u'|^2-\lambda|u|^2)\)。使用区间正弦展开所得 \(\lVert u\rVert_2\le\lVert u'\rVert_2\)，分别讨论 \(\lambda<1\)、\(\lambda=1\)、\(\lambda>1\) 的极小元与强制性。
\end{exercise}
\begin{solution}
若 \(\lambda<1\)，当 \(0\le\lambda<1\) 时 \(J_\lambda\ge(1-\lambda)\lVert u'\rVert_2^2/2\)，当 \(\lambda<0\) 时更有 \(J_\lambda\ge\lVert u'\rVert_2^2/2\)，故强制且唯一极小元为零。若 \(\lambda=1\)，正弦展开给出 \(2J_1=\frac\pi2\sum_{n\ge1}(n^2-1)|a_n|^2\)，所以极小元恰为 \(c\sin x\)，泛函不强制。若 \(\lambda>1\)，沿 \(u=c\sin x\) 有 \(J_\lambda=\pi(1-\lambda)c^2/4\to-\infty\)，故没有有限极小值。展开公式先对有限正弦和成立，再由 \(H_0^1\) 逼近推广。
\end{solution}

% 自拟：非对称双线性型的有限维反例。
\begin{exercise}\label{b2:ex:14-skew-part}
令 \(A=\left(\begin{smallmatrix}2&-3\\3&2\end{smallmatrix}\right)\)、\(f=(1,0)\)。分别求 \(Au=f\) 的解及 \(\tfrac12u^{\mathsf T}Au-f\cdot u\) 的极小元。证明 \((Au)\cdot u\ge2|u|^2\)，并解释这一强制性为何不足以使两个问题相同。
\end{exercise}
\begin{solution}
\(A^{-1}=13^{-1}\left(\begin{smallmatrix}2&3\\-3&2\end{smallmatrix}\right)\)，故方程解为 \((2/13,-3/13)\)。由于 \(u^{\mathsf T}Au=2|u|^2\)，泛函是 \(|u|^2-u_1\)，唯一极小元为 \((1/2,0)\)。强制性等式确实成立，但二次型只保留 \((A+A^{\mathsf T})/2=2I\)；驻点方程为 \(2u=f\)，反对称部分没有进入能量。
\end{solution}

% 自拟：障碍约束产生变分不等式。
\begin{exercise}\label{b2:ex:14-obstacle}
设 \(K\) 是 \(H_0^1(\Omega)\) 的非空闭凸集，\(f\in L^2(\Omega)\)，\(\Omega\) 有界。证明 Dirichlet 外力能量在 \(K\) 上有唯一极小元 \(u\)，并证明它等价于
\[\int_\Omega\nabla u\cdot\nabla(v-u)\ge\int_\Omega f(v-u)\qquad(v\in K).\]
\end{exercise}
\begin{solution}
Poincaré 不等式给出强制性；闭凸集弱闭，能量凸且范数连续，故弱下半连续。第一卷定理 9.5.9 给出存在性，梯度平方的严格凸性给出唯一性。若 \(u\) 极小，对 \(u+t(v-u)\in K\)、\(0<t\le1\) 展开能量，除以 \(t\) 令 \(t\downarrow0\)，即得不等式。反之，取 \(h=v-u\) 于精确二次展开，线性项由不等式非负，平方项也非负，故 \(E(v)\ge E(u)\)。凸约束只提供单侧可行线段，因而得到不等式。
\end{solution}

% 自拟：受迫阻尼波的能量估计。
\begin{exercise}\label{b2:ex:14-damped-wave}
设光滑零边界函数满足 \(u_{tt}-c^2\Delta u+\gamma u_t=f\)，其中 \(\gamma>0\)。证明对 \(0\le t\le T\)，
\[\mathcal E(t)+\frac\gamma2\int_0^t\lVert u_s\rVert_2^2\,\mathrm ds
 \le\mathcal E(0)+\frac1{2\gamma}\int_0^t\lVert f(s)\rVert_2^2\,\mathrm ds.\]
说明该估计怎样给出相同初边值数据的唯一性。
\end{exercise}
\begin{solution}
乘方程以 \(u_t\) 并分部积分，得 \(\mathcal E'+\gamma\lVert u_t\rVert_2^2=\int fu_t\)。由 Young 不等式，右端不超过 \(\lVert f\rVert_2^2/(2\gamma)+\gamma\lVert u_t\rVert_2^2/2\)。移项并对时间积分，得到所需估计。对两个解之差应用，右端外力与初始能量都为零，故差的时间导数和梯度均为零；零边界或零初位移固定剩余常数，差恒为零。
\end{solution}
